\section{Conclusions and Future Work}
\label{sec:conclusions}
In this paper, we present the first large-scale study of \gls{DMARC} adoption and reporting from a \gls{ccTLD} perspective. We cover 301 \glspl{ccTLD} and more than 19.2M domains, including 40k governmental domains.
Our findings show European countries leading in \gls{DMARC} adoption, but with varying levels of policy stringency, where Italy fares the worst (8.1\%). We find that government domains from India mostly use stringent \gls{DMARC} policies (94.3\%).
Our analysis of \gls{DMARC} reporting revealed indications of dependencies on foreign companies, raising concerns over digital sovereignty and data governance.
From our analysis, guidelines from national authorities and registries, financial incentives, compliance requirements, and support from hosting and email providers play a crucial role in shaping DMARC adoption and reporting practices.
As future work, we encourage the community to investigate \gls{DMARC} usage on domains tied to national infrastructure, such as those in the energy, healthcare, transport, and finance sectors.
